\documentclass{amsart}
\usepackage{hyperref}
\usepackage{amsmath}
\usepackage{amsfonts}
\usepackage{amsthm}
\usepackage{amssymb}
\usepackage{mathtools}
\usepackage{tikz}
\usepackage{tikz-cd}
\usepackage{algorithm}
\usepackage{algpseudocode}
\usepackage{mathrsfs}
\usepackage{stmaryrd}
\usepackage{stackengine}
\usepackage[margin=1in]{geometry}

\newtheorem{lemma}{Lemma}[section]
\newtheorem{corollary}[lemma]{Corollary}
\newtheorem{proposition}[lemma]{Proposition}
\newtheorem{theorem}{Theorem}[section]

\theoremstyle{definition}
\newtheorem{definition}[lemma]{Definition}
\newtheorem{example}[lemma]{Example}

\theoremstyle{remark}
\newtheorem*{remark}{Remark}

\numberwithin{equation}{section}

\DeclareMathOperator{\sch}{\mathfrak{S}}
\newcommand{\BRC}{\mathcal{BRC}}
\newcommand{\RC}{\mathcal{RC}}
\newcommand{\grass}{\mathcal{G}}
\newcommand{\wof}[1]{w_{#1}}
\newcommand{\wtt}{\mathtt{wt}}
\newcommand{\hr}{\mathtt{ht}}
\newcommand{\trm}{\mathtt{trim}}
\newcommand{\zeromap}{\mathcal{Z}}
\newcommand{\rcdd}{\mathfrak{F}}
\newcommand{\shup}{{\uparrow}}
\newcommand{\downvar}[1]{\stackrel{#1}{\searrow}}
\newcommand{\maxd}{\mathfrak{m}}
\newcommand{\code}{\mathfrak{c}}
\newcommand{\ins}[2]{{#2\overset{#1}{\leadsto}}}

\title{The bounded RC factor algebra}
\author{Matthew J. Samuel}

\begin{document}

\begin{abstract}
We present a self-contained account of the bounded RC factor algebra, an associative algebra whose basis elements are ordered tensors of full-rank Grassmannian RC graphs.  The construction rests on three pillars: the ring of bounded RC graphs, equipped with a product defined via clipping and trimming; the Demazure crystal structure on RC graphs together with a family of crystal divided difference operators; and the squash product, which realizes the tensor product of crystals on the level of pipe dreams.  After developing these ingredients, we define the algebra by means of confluent rewriting relations on tensors of Grassmannian factors and describe how to expand Schubert polynomials in this basis via the crystal elementary monomial representation.
\end{abstract}

\maketitle

\section{Introduction}\label{section:intro}

Schubert polynomials, introduced by Lascoux and Sch\"utzenberger, represent cohomology classes of Schubert varieties in the flag manifold and form a distinguished basis of the polynomial ring.  Their structure constants, the Littlewood--Richardson coefficients for Schubert calculus, remain the subject of intense combinatorial investigation.  RC graphs (also known as pipe dreams or reduced compatible sequences) provide a natural combinatorial indexing of Schubert monomials, and the additional structure one can impose on them has been a source of considerable progress.

The goal of this paper is to construct an associative algebra, here called the bounded RC factor algebra, whose elements are formal linear combinations of ordered tensors of full-rank Grassmannian RC graphs.  We write a typical basis element as $G_1 \otimes G_2 \otimes \cdots \otimes G_k$ where each $G_i$ is a Grassmannian pipe dream whose unique descent is equal to the number of rows in $G_i$.  The algebra structure is defined by concatenation of tensors followed by a confluent normalization procedure.  The resulting algebra maps surjectively onto the ring of bounded RC graphs via the squash product, and the Schubert element $\sch_w$ can be expressed as a finite linear combination of these tensor basis elements through what we call the crystal elementary monomial (CEM) representation.

We proceed as follows.  In Section~\ref{section:bounded-rc} we recall the definition of RC graphs and bounded RC graphs, and we define the ring structure on the free abelian group they span.  Section~\ref{section:crystals} develops the Demazure crystal structure and the crystal divided difference operators.  Section~\ref{section:zeroing} treats the zeroing transition algorithm, which is the key technical ingredient for the ring product.  Section~\ref{section:squash} introduces the squash product and the squash decomposition of an RC graph into an $S_n$-factor and a Grassmannian factor.  Section~\ref{section:algebra} defines the bounded RC factor algebra and its normalization relations.  Finally, Section~\ref{section:cem} describes the CEM representation of Schubert polynomials in this algebra.




\section{Bounded RC graphs and their ring structure}\label{section:bounded-rc}

\subsection{RC graphs}

Let $\mathbb{P} = \{1, 2, 3, \ldots\}$ denote the positive integers and let $S_\infty$ denote the group of permutations of $\mathbb{P}$ that fix all but finitely many elements, generated by the adjacent transpositions $s_i = (i, i+1)$ for $i \geq 1$.

\begin{definition}\label{definition:rcgraph}
Let $R \subseteq \mathbb{P} \times \mathbb{P}$ be a finite set.  To each such set we associate a permutation $\wof{R} \in S_\infty$ as follows.  Given a pair $(i, j)$ with $i, j > 0$, define $s(i,j) = s_{i+j-1}$.  Totally order $\mathbb{P} \times \mathbb{P}$ by declaring $(i, j) < (a, b)$ if and only if $i < a$, or $i = a$ and $j > b$.  Index the elements of $R$ as $r_1, r_2, \ldots, r_m$ in this order and set
\[
\wof{R} = s(r_1) \cdots s(r_m).
\]
If $\ell(\wof{R}) = m$, then $R$ is called an \textit{RC graph} (or a \textit{pipe dream}).
\end{definition}

Informally, an RC graph is a collection of crossing tiles placed in positions $(i, j)$ of a grid, read from top to bottom and right to left within each row, such that the resulting product of transpositions is a reduced word for the associated permutation.

\begin{definition}\label{definition:boundedrc}
A \textit{bounded RC graph} is an RC graph $R$ together with a positive integer $n \geq \max\{i : (i, j) \in R \text{ for some } j\}$ that satisfies $\maxd(\wof{R}) \leq n$, where $\maxd(w)$ denotes the largest descent of $w$.  We write $\hr(R) = n$ for the \textit{height} of the bounded RC graph and $\wtt(R)$ for its \textit{weight vector}, where $\wtt_i(R)$ counts the number of elements in row $i$.

The \textit{trim} operation removes the top row: $\trm(R)$ is the bounded RC graph of height $\hr(R) - 1$ with underlying set $\{(i-1, j-1) : (i, j) \in R,\, i \geq 2\}$.  The \textit{shift} operation acts by $\shup^m R = \{(i + m, j) : (i, j) \in R\}$.
\end{definition}

Let $\BRC$ denote the free abelian group spanned by all bounded RC graphs.  It admits a grading $\BRC = \bigoplus_{n \geq 0} \BRC_n$ where $\BRC_n$ is spanned by those $R$ with $\hr(R) = n$.

\subsection{The clipping and trimming operations}

The ring product on $\BRC$ requires the ability to break a bounded RC graph into a ``top part'' and a ``bottom part'' and to reassemble it.  This is accomplished by two operations.  Given a bounded RC graph $R$ with $\hr(R) = n$ and an integer $0 \leq p \leq n$, we define:

\begin{definition}\label{definition:clip}
The \textit{clip} of $R$ to height $p$, denoted $\mathtt{clip}^p(R)$, is the bounded RC graph of height $p$ obtained by taking the restriction of $R$ to its first $p$ rows and applying the zeroing transition (described in Section~\ref{section:zeroing}) to remove all descents larger than $p$.

The \textit{trim} to depth $p$, denoted $\trm^p(R)$, is the bounded RC graph of height $n - p$ given by $\trm^p(R) = \{(i - p, j) : (i, j) \in R,\, i > p\}$.
\end{definition}

These two operations satisfy the following associativity properties, which are crucial for the ring product.

\begin{lemma}\label{lemma:clip-assoc}
For a bounded RC graph $R$ and integers $0 \leq p \leq p + q \leq \hr(R)$,
\begin{align}
\mathtt{clip}^p(\mathtt{clip}^{p+q}(R)) &= \mathtt{clip}^p(R), \label{eq:clip1}\\
\trm^p(\mathtt{clip}^{p+q}(R)) &= \mathtt{clip}^q(\trm^p(R)), \label{eq:clip2}\\
\trm^{p+q}(R) &= \trm^q(\trm^p(R)). \label{eq:clip3}
\end{align}
\end{lemma}

The key structural result is that the pair $(\mathtt{clip}^p(R),\, \trm^p(R))$ together with $\wof{R}$ uniquely determines $R$.

\begin{theorem}\label{theorem:injectivity}
For each integer $0 \leq p \leq n$, the map $\BRC_n \to \BRC_p \times \BRC_{n-p} \times S_\infty$ defined by $R \mapsto (\mathtt{clip}^p(R),\, \trm^p(R),\, \wof{R})$ is injective.
\end{theorem}

\subsection{The ring product}

\begin{definition}\label{definition:diamond}
The \textit{diamond product} of two bounded RC graphs $R_1$ with $\hr(R_1) = m$ and $R_2$ with $\hr(R_2) = n$ is
\[
R_1 \diamond R_2 = \sum_{\substack{R' \in \BRC_{m+n} \\ \mathtt{clip}^m(R') = R_1 \\ \trm^m(R') = R_2}} R'.
\]
\end{definition}

\begin{theorem}\label{theorem:ring}
The product $\diamond$ turns $\BRC$ into an associative ring.
\end{theorem}

\begin{proof}
Associativity follows from the identities in Lemma~\ref{lemma:clip-assoc}.  If $R_w$ is any term in the expansion of $(R_1 \diamond R_2) \diamond R_3$ with $\hr(R_w) = p + q + r$, then equation~\eqref{eq:clip1} gives $\mathtt{clip}^p(\mathtt{clip}^{p+q}(R_w)) = \mathtt{clip}^p(R_w) = R_1$, equation~\eqref{eq:clip2} gives $\trm^p(\mathtt{clip}^{p+q}(R_w)) = \mathtt{clip}^q(\trm^p(R_w)) = R_2$, and equation~\eqref{eq:clip3} gives $\trm^{p+q}(R_w) = R_3$.  Each term appears with multiplicity one, so the two bracketings produce identical sums.
\end{proof}

The ring $\BRC$ admits a natural map to the divided difference algebra that is compatible with this product, but we will not need this here.




\section{The Demazure crystal structure}\label{section:crystals}

\subsection{Crystal operators on RC graphs}

Assaf and Schilling defined a Demazure crystal structure on the set $\RC(w)$ of all RC graphs for a fixed permutation $w$.  For each $i \geq 1$, lowering operators $f_i : \RC(w) \to \RC(w) \cup \{0\}$ and raising operators $e_i : \RC(w) \to \RC(w) \cup \{0\}$ are defined via a pairing algorithm between rows $i$ and $i+1$.

Given an RC graph $R$, one pairs elements of row $i$ with elements of row $i+1$ greedily from right to left: the element $(i, j)$ is paired with the element $(i+1, k)$ where $k$ is the smallest value such that $k \geq j$ and $(i+1, k)$ is not yet paired.  Unpaired elements in row $i$ form the set $R_i$; unpaired elements in row $i+1$ form the set $L_i$.  The lowering operator $f_i$ removes the leftmost element of $R_i$ and inserts an element in the corresponding position in row $i+1$, while $e_i$ removes the rightmost element of $L_i$ and inserts one in row $i$.  If the relevant set is empty, the operator returns $0$.

These operators satisfy the Stembridge axioms for a Demazure crystal, and the connected components under crystal operators decompose $\RC(w)$ into Demazure crystal components.  Each component has a unique \textit{highest weight} element $R^{\mathrm{hw}}$ (characterized by $e_i(R^{\mathrm{hw}}) = 0$ for all $i$) and a unique \textit{lowest weight} element (characterized by $f_i = 0$ for all $i$).

\subsection{Crystal divided differences}

\begin{definition}\label{definition:crystaldivdiff}
Let $R$ be an RC graph and let $i > 0$.  Define $\rcdd^i R$ as follows.  Set $\rcdd^i R = 0$ unless $s_i$ is a right descent of $\wof{R}$ and there exists $(i, j) \in R$ with $\mathtt{rt}_R(i, j) = (i, i+1)$.  When these conditions hold, let $R' = R \setminus \{(i, j)\}$.  If $e_i(R') \neq 0$, set $\rcdd^i R = 0$.  Otherwise,
\[
\rcdd^i R = \sum_{p=0}^{\varphi_i(R')} f_i^p R',
\]
where $\varphi_i(R')$ is the string length in the $i$-direction of the crystal.
\end{definition}

The crystal divided difference operators satisfy the nilCoxeter relations: $(\rcdd^i)^2 = 0$, $\rcdd^i \rcdd^j = \rcdd^j \rcdd^i$ for $|i - j| > 1$, and $\rcdd^i \rcdd^{i+1} \rcdd^i = \rcdd^{i+1} \rcdd^i \rcdd^{i+1}$.

\begin{theorem}\label{theorem:divdiff-schubert}
If $i$ is a right descent of $w$, then $\rcdd^i \mathcal{S}_w(n) = \mathcal{S}_{ws_i}(n)$, where $\mathcal{S}_w(n) = \sum_{\wof{R} = w,\, \hr(R) = n} R$ is the Schubert element.  If $i$ is not a right descent, then $\rcdd^i \mathcal{S}_w(n) = 0$.
\end{theorem}

A consequence is that the \textit{principal RC graph} $R^0(w)$, defined as the unique RC graph with $(i, j) \in R^0(w)$ if and only if $1 \leq j \leq \code_i(w)$, is the only RC graph from which one can successively apply crystal divided differences along any reduced word for $w$ without obtaining zero.




\section{The zeroing transition}\label{section:zeroing}

The zeroing transition algorithm is the mechanism by which one removes an empty row from the bottom of a bounded RC graph, adjusting the crossings so that the resulting graph is a valid bounded RC graph of smaller height.  It underlies both the clipping operation and the diamond product.

\begin{definition}\label{definition:zeroing}
Let $R$ be a bounded RC graph with $\hr(R) = n$ and the last row of $R$ empty.  The \textit{zeroing map} $\zeromap(R)$ produces a bounded RC graph $R'$ with $\hr(R') = n - 1$, $|R'| = |R|$, and $\wof{R} \downvar{n} \wof{R'}$ (that is, $\wof{R'}$ lies in the image of $\wof{R}$ under the parabolic projection at position $n$).  The procedure is described in Algorithm~\ref{algorithm:zeroing}.
\end{definition}

\begin{algorithm}
\caption{The zeroing map $\zeromap(R)$}
\label{algorithm:zeroing}
\begin{algorithmic}[1]
\State Let $n = \hr(R)$ and suppose row $n$ of $R$ is empty.
\If{removing row $n$ from $R$ yields a valid bounded RC graph}
  \State \Return $R$ with row $n$ removed.
\Else
  \State Set $R_0 \gets R$ and find the maximal sequence $(i_1, j_1), \ldots, (i_p, j_p)$ in $R$ such that $\mathtt{rt}_{R_{m-1}}(i_m, j_m) = (n + m - 1, n + m)$, where $R_m = R_{m-1} \setminus \{(i_m, j_m)\}$.
  \State Set $R^- = R \setminus \{(i_1, j_1), \ldots, (i_p, j_p)\}$ and $I = (i_1, \ldots, i_p)$.
  \State Form $D = R^- \cup \{(n, 1), (n, 2), \ldots, (n, p)\}$.
  \State Apply pseudo-Pieri insertion: $D' = \ins{n-1}{I} D$.
  \State \Return $R' = \{(i, j) \in D' : i < n\}$ with $\hr(R') = n - 1$.
\EndIf
\end{algorithmic}
\end{algorithm}

The zeroing map is a weight-preserving bijection from the set of $n$-row bounded RC graphs $R$ with empty row $n$ to bounded RC graphs of height $n - 1$ with appropriately related permutations.  The clipping operation $\mathtt{clip}^p(R)$ is defined by restricting $R$ to its first $p$ rows and then iteratively applying $\zeromap$ until all descents above $p$ have been removed.

The correctness of this algorithm relies on a reconstruction lemma: given the ``stripped'' graph $\overline{R}$ and the row sequence $\mathtt{strip}(R)$ obtained by removing crossings that carry descents at or above $n$, one recovers $R$ by pseudo-Pieri insertion.  The zeroing map also coincides, under the Edelman--Greene correspondence, with the Little bump operation, which preserves the recording tableau and hence the crystal structure.




\section{The squash product and squash decomposition}\label{section:squash}

\subsection{Grassmannian RC graphs}

A permutation $w$ is called $n$-\textit{Grassmannian} if its only descent is at position $n$.  The corresponding Schubert polynomial $\sch_w(x_1, \ldots, x_n)$ is a Schur polynomial in $n$ variables.  An RC graph for such a permutation, when viewed as a bounded RC graph of height $n$, is called a Grassmannian RC graph.

We say that a Grassmannian RC graph is \textit{full rank} (or \textit{full}) if its unique descent equals $n$, the number of rows.

\begin{definition}\label{definition:grassring}
For each $n$, let $\grass_n$ denote the free abelian group generated by the RC graphs $G$ with $\hr(G) = n$ such that $\wof{G}$ is $n$-Grassmannian.  Under the squash product (defined below), $\grass_n$ is a ring isomorphic to the ring of tableaux under reverse Knuth insertion.
\end{definition}

\subsection{The squash product}

Given two RC graphs $R_1$ and $R_2$, their \textit{disjoint union} $R_1 \star R_2$ is formed by shifting the columns of $R_2$ far enough to the right that no column overlaps with $R_1$, then taking the union.  Formally, choosing $N$ larger than $i + j$ for all $(i, j) \in R_2$, one sets
\[
R_1 \star R_2 = R_1 \cup \{(i, j + N) : (i, j) \in R_2\}.
\]

\begin{definition}\label{definition:squash}
The \textit{squash product} of two bounded RC graphs $R_1$ and $R_2$ with $\hr(R_1) = \hr(R_2) = n$ is
\[
R_1 \boxtimes R_2 = \mathtt{clip}^n(R_1 \star R_2).
\]
\end{definition}

The squash product is associative and preserves the crystal structure in the following sense.

\begin{theorem} \label{theorem:starproductcrystal}
  Suppose $(R_1,n)$ and $(R_2,n)$ are bounded RC graphs. Then we have that there is a homomorphism of crystal graphs
  $$\mathscr{C}(u,n)\otimes\mathscr{C}(v,n)\to \mathscr{C}(u) \star_n \mathscr{C}(v) $$
   given by $R_1\otimes R_2 \mapsto R_1\star R_2$. If $R_2$ is $n$-Grassmannian, then this is a weight-preserving isomorphism of crystal graphs.
\end{theorem}
\begin{proof}
Clearly the map is a weight preserving bijection. We need only show that it commutes with the crystal operators. Let $R_1\in \mathscr{C}(u,n)$ and $R_2\in \mathscr{C}(v,n)$, and let $R=R_1\star R_2$. Consider $f_i(R_1\otimes R_2)$. If $\varphi_i(R_1)>\varepsilon_i(R_2)$, then the unpaired elements in row $i$ in $R_1$ exhaust the unpaired elements in row $i+1$ in $R_2$. Therefore, only elements of $R_1$ are affected by $f_i$, as in the definition of the tensor product. If instead $\varphi_i(R_1)\leq \varepsilon_i(R_2)$, then all unpaired elements of $R_1$ in row $i$ are paired with elements of $R_2$ in row $i+1$, and so only elements of $R_2$ are affected by $f_i$. The same reasoning applies to $e_i$.
\end{proof}


\begin{theorem}\label{theorem:squash-crystal}
Let $u, v \in S_\infty$ have height $n$, with $v$ being $n$-Grassmannian.  The map $A \otimes B \mapsto A \boxtimes B$ is a weight-preserving crystal isomorphism from $\RC(u) \otimes \RC(v)$ to $\bigoplus_{w : c_{u,v}^w > 0} \RC(w)^{\oplus c_{u,v}^w}$.
\end{theorem}

The squash product endows $\grass_n$ with a ring structure and $\BRC_n$ with the structure of a right $\grass_n$-module.

\begin{definition}
Define 
$$\mathrm{Schub}_{\RC}(w, n) = \sum_{\wof{R} = w,\, \hr(R) = n} R$$ 
to be the Schubert element corresponding to $w$ in $\BRC_n$. Also define
$$\mathrm{DSchub}_{\RC}(w, n;y) = \sum_{uv = w, \ell(u) + \ell(v)=\ell(v)} \mathrm{Schub}_{\RC}(v, n)\sch_{u^{-1}}(y)$$ 
Also define
$$\mathrm{DElem}_{p,k}(n;y) = \sum_{\sch_v(x)=e_{p,k}(x)} \sch_{u^{-1}}(y)\mathrm{DSchub}_{\RC}(v, n; y)$$
\end{definition}

The following result is key to the construction of the bounded RC factor algebra.

\begin{theorem}
Let $\lambda$ be a partition, and let $\mu$ be the corresponding dominant permutation such that $\code^*(\mu)=\lambda$. Let $\overline{\lambda}=(\lambda_{\ell(\lambda)}, \lambda_{\ell(\lambda)-1},\ldots, \lambda_1)$. Then
$$\bigotimes_{i = 1}^{\ell(\lambda)}\mathrm{DElem}_{\overline{\lambda}_{i}, \overline{\lambda}_{i}}(\overline{\lambda}_{i};y_{\ell(\lambda) + 1 - i})$$
$$=\sum_{\ell(u\mu^{-1}) = \ell(\mu) - \ell(u)} \mathrm{Schub}_{\RC}(u, n)\sch_{u\mu^{-1}}(y)$$
\end{theorem}

\subsection{Squash decomposition}

The squash decomposition reverses the squash product for a single RC graph, factoring it into an $S_n$-piece and a Grassmannian piece.

\begin{lemma}
Suppose $R \in \BRC_n$ is a bounded RC graph of height $n$ that is a highest weight element. Then for each $m$ such that there is a pair $(R_\ast, R^\ast)$ such that $\wof{R_\ast} \in S_n$ and $R^\ast \in \grass_n$ with $|R^*|=m$ satisfying
\[
R_\ast \boxtimes R^\ast = R.
\]
we have that $R_\ast$ and $R^\ast$ are unique.
\end{lemma}

\begin{proof}
% It is shown in \cite{notes} that the polynomial ring $\mathbb{Z}[x_1,\ldots,x_n]$ is a free module over the ring $\Lambda_n$ of symmetric polynomials with basis given by the Schubert polynomials $\sch_w$ with $w\in S_n$.  This is equivalent to saying that the set of polynomials $\sch_w(x)s_\lambda(x)$ with $w\in S_n$ and $s_\lambda$ a Schur polynomial in $n$ variables is a $\mathbb{Z}$-basis for $\mathbb{Z}[x_1,\ldots,x_n]$.  By Theorem \ref{theorem:squash-crystal}, this theorem also holds at the level of RC graphs. Broadly, this means that the tensors
% $$\sum_{R\in\RC_n(w)} R \otimes \sum_{T\in \RC(\lambda)} T$$
% map to a basis of a free $\Lambda_n$-submodule of $\BRC_n$ under the squash product. 
We proceed by induction on $m$. The base case $m=0$ is clear, as in that case $\wof{R}$ is in $S_n$. Assuming the result for $m-1$, let $R$ be such that there exists an $R^\ast$ with $|R^\ast|=m$ satisfying the equation for some $R_*$. Let $R_2$ and $R^2$ be any other choice of decomposition such that $|R^2|=m$ and $R_2\in \RC_n(S_n)$. Since $R$ is a highest weight element, $R_2$ must be a highest weight element of $\RC_n(S_n)$. By the induction hypothesis, for any RC graph $A$ with $|A| = |R|-1$ such that there exists a degree $1$ RC graph $M\in\grass_n$ for which
$R = A \boxtimes M,$ we have that $A$ admits a unique squash decomposition $A = A_\ast \boxtimes A^\ast$ with $|A^\ast|=m-1$, since $A$ must itself be a highest weight. By invertibility of Schensted insertion, it follows that there is a bijection between the set of such $A$ and the set of degree $1$ RC graphs in $\grass_n$ that can be factored out of $R$. This implies uniqueness of $R^*$  by induction. For such an $R^*$ there is a corresponding $A_*$ such that $R = A_* \boxtimes R^*$. It follows from the fact that
$$R=A_*\boxtimes A^*\boxtimes M$$
and the fact that $\boxtimes$ is a right module action that $A_*=R_2$, and we are done.
\end{proof}

\begin{proposition}\label{proposition:unique-factorization}
Given a bounded RC graph $R$ of height $n$ and a weak composition $\alpha$ of length $n$ whose sum is $|R|$, there is at most one factorization of the form
$$E_1\boxtimes E_2\boxtimes \cdots \boxtimes E_{n-1}\boxtimes G = R$$
where each $E_i$ is an RC graph for an elementary symmetric polynomial in $i$ variables, $G$ is a full-rank $n$-Grassmannian RC graph, and $\wtt(E_i) = \alpha_i$ for all $i<n$.
\end{proposition}
\begin{proof}

\end{proof}

% Let $\RC^{\alpha}_{n-1}(S_n)$ denote the set of RC graphs $R$ of height $n-1$ such that $\wof{R}\in S_n$ and $\wtt^{\mathrm{col}}(R) = \alpha$. 
% \begin{proposition}
% We have that
% $$E_{\alpha_1,1}\boxtimes E_{\alpha_2,2}\boxtimes \cdots \boxtimes E_{\alpha_{n-1},n-1} = \sum_{R\in \RC^{\alpha}_{n-1}(S_n)}R$$
% \end{proposition}

\begin{theorem}
Fix a positive integer $n$. Let $\mu$ be a dominant permutation such that 
$$\code^*(\mu)=(n^p, n-1, n-2, \ldots, 1)$$
for some $p\geq 0$. Then we have
\begin{multline*}
\sum_{\alpha}x_1^{n + p - \alpha_1}x_2^{n + p - 1 - \alpha_2}\cdots x_{n+p-1}^{1 - \alpha_{n+p-1}}
E_{\alpha_{n+p-1},1}\boxtimes  \cdots \boxtimes E_{\alpha_{p+1},n -1}\boxtimes E_{\alpha_p,n}\boxtimes \cdots \boxtimes E_{\alpha_1,n} \\
= \sum_{\substack{R\in\RC(w),\, w:\\[2pt] \ell(w\mu^{-1})=\ell(\mu)-\ell(w)}}R\,\sch_{w\mu^{-1}}(x)
\end{multline*}
where the sum on the left is over all weak compositions $\alpha$ of length $n+p-1$ such that $\alpha_i\leq n$ for $1\leq i\leq p$ and $\alpha_i\leq n - (i-p)$ for $i>p$. 
\end{theorem}
\begin{proof}
First assume $p=0$. We prove the result by induction on $n$. For $n=1$, this is a triviality. For larger $n$, assuming the result for $n-1$, we have that
\begin{multline*}\sum_{\alpha}x_1^{n - \alpha_1}x_2^{n - 1 - \alpha_2}\cdots x_{n-1}^{1 - \alpha_{n-1}}
E_{\alpha_{n-1},1}\boxtimes  \cdots \boxtimes E_{\alpha_1,n} \\
= \sum_{\alpha}x_1^{n - \alpha_1}x_2^{n - 1 - \alpha_2}\cdots x_{n-1}^{1 - \alpha_{n-1}}
\left(\sum_{R\in \RC^{\alpha}_{n-1}(S_{n-1})}R\right)\boxtimes E_{\alpha_1,n} \\
= \sum_{\alpha}x_1^{n - \alpha_1}x_2^{n - 1 - \alpha_2}\cdots x_{n-1}^{1 - \alpha_{n-1}}
\left(\sum_{\substack{R\in \RC_{n-1}(w),\, w\in S_{n-1}:\\[2pt] \ell(w\mu^{-1})=\ell(\mu)-\ell(w)}}R\right)\boxtimes E_{\alpha_1,n} \\
= \sum_{\alpha}x_1^{n - \alpha_1}x_2^{n - 1 - \alpha_2}\cdots x_{n-1}^{1 - \alpha_{n-1}}
\left(\sum_{\substack{R\in \RC_{n-1}(w),\, w\in S_{n-1}:\\[2pt] \ell(w\mu^{-1})=\ell(\mu)-\ell(w)}}R\boxtimes E_{\alpha_1,n}\right) \\
= \sum_{\alpha}x_1^{n - \alpha_1}x_2^{n - 1 - \alpha_2}\cdots x_{n-1}^{1 - \alpha_{n-1}}
\left(\sum_{\substack{R\in \RC_{n-1}(w),\, w\in S_{n-1}:\\[2pt] \ell(w\mu^{-1})=\ell(\mu)-\ell(w)}}\sum_{\substack{R'\in \RC(w'),\, w':\\[2pt] \ell(w'\mu^{-1})=\ell(\mu)-\ell(w')}}R'\right) \\
= \sum_{\substack{R\in\RC(w),\, w:\\[2pt] \ell(w\mu^{-1})=\ell(\mu)-\ell(w)}}R\,\sch_{w\mu^{-1}}(x)
\end{multline*}
where the third equality uses Theorem~\ref{theorem:starproductcrystal} and the fourth uses Theorem~\ref{theorem:squash-crystal}.

For the case $p>0$, the factors on the left-hand side separate into two groups: those of sizes $1$ through $n-1$ (namely $E_{\alpha_{n+p-1},1}, \ldots, E_{\alpha_{p+1},n-1}$) and $p$ factors of size $n$ (namely $E_{\alpha_p,n}, \ldots, E_{\alpha_1,n}$).  By associativity of $\boxtimes$, we may write each term as
\[
\left(E_{\alpha_{n+p-1},1}\boxtimes \cdots \boxtimes E_{\alpha_{p+1},n-1}\right) \boxtimes \left(E_{\alpha_p,n}\boxtimes \cdots \boxtimes E_{\alpha_1,n}\right).
\]
The factors of sizes $1, \ldots, n-1$, summed over $(\alpha_{p+1}, \ldots, \alpha_{n+p-1})$ with the monomials $x_{p+1}^{n - \alpha_{p+1}} \cdots x_{n+p-1}^{1 - \alpha_{n+p-1}}$, constitute an instance of the $p = 0$ result (with variables shifted by $p$) applied to the dominant permutation $\mu'$ with $\code^*(\mu') = (n-1, n-2, \ldots, 1)$.  The squash product $E_{\alpha_p,n} \boxtimes \cdots \boxtimes E_{\alpha_1,n}$ yields a single $n$-Grassmannian RC graph $G$.  By Proposition~\ref{proposition:unique-factorization}, each bounded RC graph $R$ of height $n$ admits at most one factorization $E_1 \boxtimes \cdots \boxtimes E_{n-1} \boxtimes G = R$ with prescribed weights, so the identification of RC graphs between the two sides is injective considering each weight.  The result then follows from Theorems~\ref{theorem:starproductcrystal} and~\ref{theorem:squash-crystal}, which provide the surjection and the correct Schubert polynomial coefficients.
\end{proof}

\section{The bounded RC factor algebra}\label{section:algebra}

\subsection{Basis elements}

Having developed the squash product, the squash decomposition, and the Demazure crystal structure, we are now in a position to define the algebra.

\begin{definition}\label{definition:brfa}
A \textit{factor key} of \textit{size} $n$ is an ordered tuple $(G_1, G_2, \ldots, G_k)$ of full-rank Grassmannian RC graphs, together with a positive integer $n$ satisfying $n \geq \max_i \hr(G_i)$, where each $G_i$ is an $n_i$-Grassmannian RC graph with $\hr(G_i) = n_i$ and $n_i$ equal to the unique descent of $\wof{G_i}$ (the full-rank condition).  The \textit{bounded RC factor algebra} $\mathcal{E}$ is the free $\mathbb{Z}$-module with basis the set of all factor keys in normal form (defined below), equipped with the product given by concatenation followed by normalization.
\end{definition}

\subsection{Normalization relations}

The normalization of a factor key is a confluent rewriting system governed by three types of rules, applied iteratively until none applies.

\begin{enumerate}
\item \textit{Sorting.}  If two adjacent factors satisfy $\hr(G_i) > \hr(G_{i+1})$, they are transposed.  Factors of distinct size commute, and after sorting, the sizes weakly increase from left to right.

\item \textit{Merging.}  If two adjacent factors $G_i$ and $G_{i+1}$ satisfy $\hr(G_i) = \hr(G_{i+1})$, they are replaced by the single RC graph $G_i \boxtimes G_{i+1}$.  This merged graph is again Grassmannian and has the same number of rows.

\item \textit{Decomposition.}  If a factor $G_i$ represents a Grassmannian permutation $u$ such that $u\notin S_{\hr(G_i)}$, then $G_i$ can be resized to a larger height and decomposed via the squash decomposition into a pair $(G_i', G_i'')$ of full-rank Grassmannian factors, with $\hr(G_i')=\hr(G_i)$ and $\hr(G_i'')=\hr(G_i)+1$, such that $G_i' \boxtimes G_i'' = G_i$.  The factor key is then rewritten by replacing $G_i$ with the pair $(G_i', G_i'')$.
\end{enumerate}

In addition, any identity factor (an RC graph for the identity permutation) is dropped from the tensor.

\begin{proposition}\label{proposition:confluence}
For each $n$, every factor key of height at most $n$ has a unique normal form $G_1\otimes \cdots \otimes G_n$ obtained by applying the above rules until none applies.
\end{proposition}

The product on $\mathcal{E}$ is then defined by
\[
(G_1 \otimes \cdots \otimes G_p) \cdot (H_1 \otimes \cdots \otimes H_q) = \mathrm{normalize}(G_1 \otimes \cdots \otimes G_p \otimes H_1 \otimes \cdots \otimes H_q),
\]
extended bilinearly to arbitrary elements.

\begin{theorem}\label{theorem:associativity}
The product on $\mathcal{E}$ is associative.
\end{theorem}

\begin{proof}
Write elements as formal linear combinations of factor keys and expand.  Associativity reduces to the statement that for factor keys $\alpha$, $\beta$, $\gamma$, the normal form of $(\alpha \otimes \beta \otimes \gamma)$ does not depend on the order in which adjacent normalizations are performed.  This follows from the confluence of the rewriting system.
\end{proof}

\subsection{The squash map}

There is a natural linear surjection $\mathcal{R} : \mathcal{E} \to \BRC$ that evaluates each factor key by iterated squash products:
\[
\mathcal{R}(G_1 \otimes G_2 \otimes \cdots \otimes G_k) = (\cdots((G_1 \boxtimes G_2) \boxtimes G_3) \cdots) \boxtimes G_k.
\]
This map is well defined on normal forms because the normalization rules preserve the squash product: merging two equal-size factors via $\boxtimes$ is precisely the squash product, sorting commuting factors preserves the overall squash, and the decomposition rule is the inverse of a squash product and therefore also preserves it.

\begin{remark}
The squash map is not a ring homomorphism.  Its kernel is not a two-sided ideal of $\mathcal{E}$.  This fact, demonstrated by explicit examples already for the Schubert polynomial $\sch_{1432}$, is one of the main reasons the factor algebra is necessary: it retains more structure than the ring of RC graphs alone.
\end{remark}



\section{The crystal elementary monomial representation}\label{section:cem}

\subsection{Elementary symmetric RC graphs}

For integers $0 \leq p \leq k$, the $k$-Grassmannian permutation with Lehmer code having a single nonzero entry $p$ in position $k$ (that is, the permutation sending $k \mapsto k + p$ and cyclically adjusting) corresponds to the elementary symmetric polynomial $e_p(x_1, \ldots, x_k)$.  Its RC graphs form a family $\{M_A^k\}$ indexed by subsets $A \subseteq [k]$ of size $p$.  Each such graph $M_A^k$ is a full-rank Grassmannian RC graph of height $k$.

The sum
\[
\mathbf{e}_p^k = \sum_{\substack{A \subseteq [k] \\ |A| = p}} M_A^k
\]
represents the elementary symmetric polynomial $e_p(x_1, \ldots, x_k)$ at the level of the algebra $\mathcal{E}$.

\subsection{CEM expansion of Schubert polynomials}

Every Schubert polynomial can be written as a polynomial in the elementary symmetric polynomials, by a recursive procedure using divided differences.  At the level of the factor algebra, this translates into the following.

\begin{definition}\label{definition:cem}
The \textit{crystal elementary monomial} (CEM) representation of a Schubert polynomial $\sch_w$ at size $n$ is the element
\[
\sch_w^{\mathrm{CEM}}(n) = \sum_{(G_1, \ldots, G_k)} c_{G_1, \ldots, G_k} \cdot (G_1 \otimes \cdots \otimes G_k) \in \mathcal{E}
\]
whose image under the squash map satisfies $\mathcal{R}(\sch_w^{\mathrm{CEM}}(n)) = \mathcal{S}_w(n)$, and where the coefficients $c_{G_1, \ldots, G_k}$ are determined by expanding $\sch_w$ in the elementary monomial basis and tracking the crystal structure.
\end{definition}

The CEM representation is computed algorithmically through the function \texttt{full\_CEM}, which expands a Schubert polynomial in terms of elementary symmetric functions via a $v$-path dictionary (implementing the Lascoux--Sch\"utzenberger transition formula), then lifts each monomial $e_{a_1}(x_1, \ldots, x_{k_1}) \cdots e_{a_m}(x_1, \ldots, x_{k_m})$ to a tensor of Grassmannian RC graphs respecting the crystal structure.

\begin{example}
Consider the permutation $w = [2, 1, 3]$ with $\sch_w = x_1$.  At size $n = 2$, the CEM representation consists of a single tensor: the Grassmannian RC graph $M_{\{1\}}^1$ of height $1$ with a single crossing at position $(1,1)$.

For $w = [3, 1, 2]$ with $\sch_w = x_1^2 + x_1 x_2 = x_1 e_1(x_1, x_2) = e_1(x_1) \cdot e_1(x_1, x_2)$ at size $n = 2$, the CEM representation is a sum of tensor products of full-rank Grassmannian factors whose squash products yield the RC graphs for $[3, 1, 2]$.
\end{example}

\subsection{Dominant permutations}

A permutation $w$ is \textit{dominant} if its Lehmer code is a partition (a weakly decreasing sequence).  For a dominant permutation $w_\lambda$ with $\code(w_\lambda) = \lambda = (\lambda_1 \geq \lambda_2 \geq \cdots \geq \lambda_m)$, the Schubert polynomial is a product of elementary symmetric polynomials:
\[
\sch_{w_\lambda} = \prod_{i = 1}^{m} e_{\lambda_i - \lambda_{i+1}}(x_1, \ldots, x_i),
\]
where $\lambda_{m+1} = 0$.  In the factor algebra, the CEM representation of a dominant permutation is therefore a single tensor monomial with coefficient $+1$, factored as a product of at most one elementary symmetric function per rank, following a pattern dictated by the differences in the partition.

For a general permutation $w$, one finds a dominant permutation $w_\lambda$ above $w$ in the left weak Bruhat order, applies crystal divided differences along a reduced word from $w_\lambda$ to $w$, and records the resulting expansion.  This is the algorithmic content of the CEM representation.

\subsection{Multiplication in the factor algebra}

Given the CEM representations of $\sch_u$ and $\sch_v$, the product $\sch_u^{\mathrm{CEM}}(n) \cdot \sch_v^{\mathrm{CEM}}(n)$ is computed in $\mathcal{E}$ by distributing and normalizing each pair of tensor monomials.  One then applies the squash map $\mathcal{R}$ to obtain an element of $\BRC_n$, and reads off the Littlewood--Richardson coefficients $c_{uv}^w$ by identifying the Schubert expansion.

The hypothesis, verified computationally for permutations up to moderate length, is that this procedure yields non-negative integer coefficients, providing a combinatorial manifestation of the Littlewood--Richardson rule for Schubert calculus.

\subsection{Formula for Schubert elements}

Let $\mu$ be a dominant permutation with $\code^*(\mu) = (\lambda_1,\ldots,\lambda_n)$ and let $u$ be a permutation such that $\ell(u\mu) = \ell(\mu) - \ell(u)$. For a positive integer $p$, we say $u$ and $w$ are \emph{$p$-anchored} if there exists a sequence of pairwise distinct positive integers $a_1,a_2,\ldots,a_k$ with $a_i\neq p$ for any $i$ such that $w=ut_{pa_1}\cdots t_{pa_k}$ and $\ell(ut_{pa_1}\cdots t_{pa_i})=\ell(u)+i$ for all $i$. We have the following formula for Schubert elements:
$$\mathrm{Schub}_{\RC}(u\mu, \lambda_1) = \sum (-1)^{\mathrm{sign}(P)}E_{\lambda_n-\ell(u_{n-1},u_n)}^{\lambda_n}\otimes \cdots \otimes E_{\lambda_2-\ell(u_1,u_2)}^{\lambda_2}\otimes E_{\lambda_1-\ell(u_0,u_1)}^{\lambda_1}$$
where $P=(u_0,\ldots,u_n)$ is a sequence of permutations with $u_0=1$ and $u_n=u$ such that $u_{i-1}$ and $u_i$ are $i$-anchored for all $i$, and $\mathrm{sign}(P)$ is the sum over all $i$ of the number of $j<i$ such that $u_{i-1}(j)\neq u_i(j)$.

\section*{Acknowledgments}

Computational verification of the constructions described here was carried out using the Python package \texttt{schubmult}.

\cite{samuelmolev}
\bibliographystyle{amsalpha}
\bibliography{schubert_bialgebra}

\end{document}
